\documentclass[a4paper]{article}
\usepackage{amsfonts}
\usepackage{amsmath}
\usepackage{amssymb}
\usepackage{graphicx}     
\usepackage{cancel}

% Command definities van frequent gebruikte commandos
\newcommand{\N}{\mathbb{N}}
\newcommand{\Q}{\mathbb{Q}}
\newcommand{\R}{\mathbb{R}}
\newcommand{\C}{\mathbb{C}}
\newcommand{\Bgsin}{\operatorname{Bgsin}}
\newcommand{\Bgcos}{\operatorname{Bgcos}}
\newcommand{\Bgtan}{\operatorname{Bgtan}}
\newcommand{\cis}{\operatorname{cis}}
\newcommand{\Limit}{\lim\limits_}
\newcommand{\Int}{\displaystyle\int}

\begin{document}
  
\noindent \large Auteur: Wout Vanderborght\\
\noindent \large Studierichting: Informatica\\
\noindent \large Oefeningenreeks 4A nr. 56.2\\

\medskip

\normalsize

$ \Int \sqrt{1+\cos 2x} \cdot dx $ \\

$ = \Int \left( \sqrt{1+\cos 2x} \cdot \dfrac{\sqrt{2}}{\sqrt{2}} \right) dx $ \\

$ = \sqrt{2} \Int \sqrt{\dfrac{1+\cos2x}{2}} \cdot dx $ \\

$ = \sqrt{2} \Int \sqrt{\cos^2 x} \cdot dx $ \\

$ = \pm \sqrt{2} \Int \cos x \cdot dx $ \\

$ = \pm \sqrt{2} \sin x + C $

\begin{thebibliography}{999}
%\bibitem{a} Referentie
\end{thebibliography}
\end{document} 
